\gddchapter{Testing sequence}
The prototype can be played with both voice and keyboard and the navigation modes can be swapped with a button press. 

\section{Baseline A - Keyboard navigation}
At first the participant will spend 10 minutes by playing the game using keyboard navigation.
This allows the player to anchor in the typical experience and serves as a comparative point in the A/B testing.

\subsection{Observation log}
\begin{table}[h]
\centering
\begin{tabular}{|p{5cm}|p{10cm}|}
\hline
\textbf{Field} & \textbf{Response} \\ \hline
\textbf{Physical comfort} How is the participant feeling & \\[1.5cm] \hline
\textbf{Speed of navigation} How quickly does the player navigate in the game & \\[1.5cm] \hline
\textbf{Immersion markers} (flow \& frustratoin) & \\[1.5cm] \hline
\end{tabular}
\end{table}

\section{Novel version B - Voice navigation}
After 10 minutes the researcher switches the navigation to voice input mode by using the keyboard. The prototype is not restarted,
instead the player is instructed to continue his playthrough using voice. The decision not to restart is placed to reduce any frustration that the player
might feel if they have to restart and repeat the actions that they just did.

\subsection{Observation log}
\begin{table}[H]
\centering
\begin{tabular}{|p{5cm}|p{10cm}|}
\hline
\textbf{Field} & \textbf{Response} \\ \hline
\textbf{Physical comfort} How is the participant feeling & \\[1.5cm] \hline
\textbf{Speed of navigation} How quickly does the player navigate in the game & \\[1.5cm] \hline
\textbf{Immersion markers} (is the player using long sentences or simple commands? What type of language do they use? Are they speaking first or 3rd person?) (flow \& frustration) & \\[1.5cm] \hline
\textbf{Natural language observation} (is the player using long sentences or simple commands? What type of language do they use? Are they speaking first or 3rd person?) & \\[1.5cm] \hline
\textbf{Linguistic guessing} (is the particpant trying to guess the possible commands by speaking to the game?) & \\[1.5cm] \hline
\textbf{Processing time user reaction} (Does the wait time break the suspension of disbelief)& \\[1.5cm] \hline
\textbf{Voice processing pipeline edge cases} (System edge cases eg. two vs to)& \\[1.5cm] \hline
\end{tabular}
\end{table}

\subsection{Sticky moments}
The moments that stood out during testing marked with timestamps.
\begin{table}[h]
\centering
\begin{tabular}{|p{3cm}|p{10cm}|}
\hline
\textbf{Timestamp} & \textbf{Log} \\ \hline
0:05 & Example timestamp \\ \hline 
 &  \\ \hline 
 &  \\ \hline 
 &  \\ \hline 
 &  \\ \hline 
 &  \\ \hline 
 &  \\ \hline 
 &  \\ \hline 

\end{tabular}
\end{table}

\section{Alignment check}
How much was the researcher participating in the gameplay experience? 
Is the participant related in any way to the researcher? Could that influence the result?

Response to be filled out.